% =====================================================================
%  qBOUNCE — the processing chain as eight stages, drawn in the manner of
%  Filter (2018) figure 6.13: what the data IS at each point and what it
%  becomes, as opposed to architecture_graphs.tex, which says which file
%  calls which. The two views are deliberately kept apart.
%
%  The thumbnail in each box is made from the real data by
%  make_flow_thumbnails.py, not drawn as an icon; stages 1 to 4 all show the
%  same 256x256 corner of the same frames, so one patch of sensor can be
%  followed from raw pixels to a labelled cluster.
%
%  Requires: tikz with arrows.meta, positioning, calc; graphicx.
% =====================================================================

\newcommand{\qbFlowBox}[5]{%
  % #1 name  #2 x  #3 y  #4 title  #5 note
  \node[stage] (#1) at (#2,#3) {};
  \node[stitle] at (#2,#3+13.5) {#4};
  \node[snote]  at (#2,#3-13.5) {#5};
}

\newcommand{\qbDataFlow}{%
\begin{tikzpicture}[x=1mm,y=1mm,
    stage/.style   = {rectangle, rounded corners=2.5pt, draw=black!55,
                      fill=black!4, minimum width=36mm, minimum height=36mm,
                      inner sep=2pt, align=center},
    stitle/.style  = {font=\small\bfseries},
    snote/.style   = {font=\scriptsize, text=black!65, align=center},
    flow/.style    = {-{Stealth[length=2.6mm,width=2.2mm]}, line width=1.1pt,
                      draw=black!75}]

  \def\dx{46}
  \def\rowa{0}
  \def\rowb{-58}

  % ── row 1: dark frames -> model -> threshold -> clusters ────────────────
  \qbFlowBox{a0}{0*\dx}{\rowa}{dark frames}{500 images, beam off\\$5472\times3672$}
  \qbFlowBox{a1}{1*\dx}{\rowa}{background model}{per-pixel $\mu$, $\sigma$\\and dead mask}
  \qbFlowBox{a2}{2*\dx}{\rowa}{threshold}{$p>\mu+5\max(\sigma,1)$\\one bit per pixel}
  \qbFlowBox{a3}{3*\dx}{\rowa}{clusters}{connected\\components}
  \foreach \i/\f in {0/s1_dark, 1/s2_model, 2/s3_threshold, 3/s4_clusters}
    \node at (\i*\dx,\rowa+1) {\includegraphics[width=17mm]{figures_flow/\f.png}};
  \draw[flow] (a0) -- (a1);
  \draw[flow] (a1) -- (a2);
  \draw[flow] (a2) -- (a3);

  % ── row 2: measurements -> labels -> classifier -> analysis ─────────────
  \qbFlowBox{b0}{0*\dx}{\rowb}{measurements}{20 columns\\per cluster}
  \qbFlowBox{b1}{1*\dx}{\rowb}{hand labels}{200 verdicts,\\three classes}
  \qbFlowBox{b2}{2*\dx}{\rowb}{classifier}{random forest,\\10 features}
  \qbFlowBox{b3}{3*\dx}{\rowb}{analysis}{rates, arrival map,\\deposited charge}
  \foreach \i/\f in {0/s5_measure, 1/s6_labels, 2/s7_classify, 3/s8_analysis}
    \node at (\i*\dx,\rowb+1) {\includegraphics[width=17mm]{figures_flow/\f.png}};
  \draw[flow] (b0) -- (b1);
  \draw[flow] (b1) -- (b2);
  \draw[flow] (b2) -- (b3);

  % the beam-on frames enter at the threshold, not at the head of the chain
  \node[stage, fill=orange!7, draw=orange!60!black,
        minimum height=11mm, minimum width=30mm] (sig) at (2*\dx, \rowa+34) {};
  \node[stitle] at (2*\dx, \rowa+36) {beam-on frames};
  \node[snote]  at (2*\dx, \rowa+31.5) {1500 images};
  \draw[flow, draw=orange!65!black] (sig) -- (a2);

  % wrap-around from the end of row 1 to the start of row 2
  \draw[flow] (a3.east) -- ++(9,0) |- (0,\rowa-27) -| (b0.north);

  % the one place a human is required
  \draw[draw=black!45, dashed, line width=0.5pt] (1*\dx, \rowb-18.5) -- (1*\dx, \rowb-21);
  \node[snote, text=black!75, anchor=north] at (1*\dx, \rowb-21)
        {\itshape the only stage that cannot be automated};

\end{tikzpicture}}
